\section{Turbulence Modeling\label{ch3:turbulenceModeling}}

The history of turbulence modeling is deeply connected to the development of
fluid dynamics and the advancements in computational power. Three primary
turbulence modeling approaches exist: \gls{DNS}, \gls{LES}, and \gls{RANS}.
\Cref{fig:turbulenceFidelity} illustrates the trade-off between computational
cost and modeling fidelity for these approaches.
\begin{figure}[H]
    \makebox[\textwidth][c]{
        \includegraphics[width=0.7\textwidth]{chapter3/turbulenceComputationalCost}
    }
    \caption{Turbulence models fidelity-cost} \label{fig:turbulenceFidelity}
\end{figure}

The choice of turbulence model depends on the specific problem and
computational resources available. At very low Reynolds numbers, where
resolving all turbulent scales is feasible, \gls{DNS} can be employed. In this
approach, the smallest resolved scale corresponds to the Kolmogorov scale,
capturing the full range of turbulence
without the need for modeling.\\

At higher Reynolds numbers, where the computational cost of resolving all
scales becomes prohibitive, \gls{LES} is used. \gls{LES} explicitly resolves
the larger, energy-containing turbulent structures while modeling the smaller,
dissipative scales. This approach strikes a balance between accuracy and
computational efficiency, making it suitable for a wide range of engineering
applications. \\

However, in most real-world applications, the Reynolds number is significantly
higher, making \gls{DNS} and even {LES} computationally expensive. As a result,
industry and large-scale engineering problems often rely on \gls{RANS} methods.
\gls{RANS} models the effects of turbulence entirely through statistical
averaging, significantly reducing computational cost while providing reasonable
accuracy for many practical applications. \\

Hybrid approaches, such as Hybrid \gls{RANS}/\gls{LES}, attempt to bridge the
gap between these methods by applying \gls{RANS} modeling in near-wall regions
while using \gls{LES} in the outer flow. This approach improves fidelity while
keeping computational costs manageable. \\

The appropriate choice of turbulence model depends on the specific objectives
of a study, the available computational resources, and the level of fidelity
required for accurate predictions \cite{Smits2009}. \\

For the \gls{DNS} case, the full \gls{NS} equations
(\cref{eq:navierStokesEquationsEins}) are resolved in both space and time
without the need for any turbulence modeling. However, the computational cost
is extremely high, making it infeasible for most real-world applications. 
\begin{subequations}\label{eq:navierStokesEquationsEins}
    \begin{align}
        &\continuityEins \label{eq:continuityEins}\\
        &\momentumEins \label{eq:momentumEins} \\
        &\energyEins \label{eq:energyEins} \\
        &\viscousStressTensor \label{eq:viscousStressTensor} \\
        &\strainRateTensor \label{eq:strainRateTensor}
    \end{align}
\end{subequations}

In $1963$, Smagorinsky \cite{Smagorinsky1963} introduced the \gls{LES} model to
simulate atmospheric air currents. The \gls{LES} framework is grounded in the
energy cascade (\cref{fig:ch2EnergyCascade}) model developed by Kolmogorov in the 
$1940$s, which itself builds upon the earlier work of Richardson. The
fundamental concept behind \gls{LES} is to explicitly compute the large-scale
motions, which are influenced flow geometry and are not universal, up to a
prescribed length scale. The smaller scales, which tend to exhibit more
universal characteristics, are modeled using a \gls{SGS} model
\cite{Pope2000}.\\

This separation between larger and smaller scales is typically achieved by
spatially filtering the velocity field using a kernel \cite{Leonard1975}. The
convolution kernel effectively splits the flow field into resolved large scales
and modeled \gls{SGS} components in an \gls{LES} framework. Developing
physically realistic \gls{SGS} models requires a comprehensive understanding of
both the physics and statistical properties of scale interactions between large
and small scales in turbulent flows. \gls{SGS} modeling is traditionally
predicated on the hypothesis that if subgrid scales exist, the flow is locally
turbulent in both space and time \cite{Garnier2009}.\\

There are two primary \gls{SGS} modeling strategies:
\begin{itemize}
    \item \textbf{Functional modeling:} This approach focuses on the energetic
        impact of subgrid scales on resolved scales. It assumes that balancing
        energy transfer between scale ranges is sufficient to describe the
        subgrid-scale effects. Functional models often neglect backscattering,
        i.e., energy transfer from smaller scales back to larger scales
        \cite{Burton2005}.
    \item \textbf{Structural modeling:} This strategy attempts to reconstruct
        subgrid-scale information directly by generating a new field from the
        filtered field (e.g., velocity, density) that better approximates the
        unfiltered field. Structural models account for backscattering because
        they aim to reconstruct small-scale structures within the
        flow \cite{Burton2005}.
\end{itemize}

Functional models commonly represent the \gls{SGS} stress tensor (Reynolds
stress tensor) using eddy-viscosity closures. However, this approach linearize
the inherently nonlinear stress tensor, leading to the loss of some physical
details. It is important to note that the famous closure problem originates
from the convective terms in the \gls{NS} equations. Additionally, functional
models are often criticized for being overly dissipative due to their nature of
transferring energy from larger to smaller scales while failing to capture
backscattering effects. In contrast, structural models incorporate
backscattering and are typically nonlinear, allowing for a more accurate
representation of small-scale flow structures.\\

If we take the \gls{NS} \cref{eq:navierStokesEquationsEins} and apply a
Reynolds decomposition \cref{eq:reynoldsDecomposition}
\begin{subequations}
    \begin{align}
        u &= \overline{u} + u^{\prime}\label{eq:reynoldsDecomposition} \\
        \overline{u} &= \lim_{T\rightarrow\infty}\frac{1}{T}\int^T_0 u\,dt \label{eq:timeAverage}
    \end{align}
\end{subequations}

We will obtain the \gls{RANS} equations give by
\cref{eq:reynoldsAverageNavierStokesEquations}.
\begin{subequations}\label{eq:reynoldsAverageNavierStokesEquations}
    \begin{align}
        &\ransContinuityEins \label{eq:ransContinuity}  \\
        &\ransMomentumEins \label{eq:ransMomentum}   \\
        &\ransEnergyEins \label{eq:ransEnergy} \\
        &\ransStressTensorEins \label{eq:ransStressTensor}
    \end{align}
\end{subequations}

The term \(\overline{u^{\prime}_i u^{\prime}_j}\) in \cref{eq:ransMomentum} is
referred to as the \textit{Reynolds stress tensor}, as shown in
\cref{eq:reynoldsStressTensor}. This is a symmetric tensor, meaning it has six
independent components. Consequently, by applying a Reynolds average to the
\gls{NS} equations, six additional unknowns emerge. However, since we only have
four equations, the system of equations remains underdetermined and is not yet
closed. This fundamental issue, arising from the application of Reynolds
averaging to the \gls{NS} equations, is known as the \textit{closure problem}.
\begin{equation}\label{eq:reynoldsStressTensor}
    \tau_{ij} = - \overline{u^{\prime}_iu^{\prime}_j} 
\end{equation}

In $1925$, Prandtl introduced the mixing length model, an analogy to the mean
free path in a gas. The model is based on the mixing length hypothesis and is
commonly referred to as an \textit{algebraic model} or a \textit{zero-equation
model} of turbulence. By definition, an $n$-equation model requires the
solution of $n$ additional differential transport equations in addition to
those governing the conservation of mass, momentum, and energy for the mean
flow \cite{Wilcox1998}. \\

In $1945$, Prandtl proposed a model in which the eddy viscosity depends on the
kinetic energy, $\kappa$, of turbulent fluctuations. He introduced a partial
differential equation to solve for $\kappa$, forming what is known as a
\textit{one-equation model} of turbulence \cite{Wilcox1998}.\\

In $1942$, Kolmogorov developed the first complete turbulence model. In
addition to a modeled equation for kinetic energy $\kappa$, he introduced a
second differential equation to solve for the dissipation rate of turbulent
energy per unit volume and time, $\omega$. This model, known as the
$\kappa$-$\omega$ model, formulates a differential equation for $\omega$
similar in structure to the equation for $\kappa$. As a result, it is
classified as a \textit{two-equation model} of turbulence \cite{Wilcox1998}.\\

Further advancements were made by Chou $1945$ and Rotta $1951$, who laid the
foundation for turbulence models that do not rely on the Boussinesq
approximation. These models involve differential equations governing the
evolution of the Reynolds stress tensor, which represents turbulent stresses.
Such models are most appropriately classified as \textit{stress transport
models} \cite{Wilcox1998}.\\

Thus, four main categories of turbulence models have evolved:
\begin{itemize}
    \item \textbf{Zero-equation models}  
    \item \textbf{One-equation models}  
    \item \textbf{Two-equation models}  
    \item \textbf{Stress transport models}  
\end{itemize}


\todo{Add these references}
\cite{DahmNotes} 
\cite{Spalart1992}








